\documentclass[11pt,oneside]{article}
\usepackage[T1]{fontenc}
\usepackage[utf8]{inputenc}
\usepackage{lmodern}
\usepackage[a4paper,margin=27mm,headheight=15pt]{geometry}
\usepackage{microtype,amsmath,amssymb,amsthm,mathtools,aliascnt}
\usepackage{booktabs,longtable,array,enumitem}
\usepackage[dvipsnames]{xcolor}
\usepackage{fancyhdr}
\usepackage{hyperref}
\hypersetup{colorlinks=true,linkcolor=MidnightBlue,citecolor=MidnightBlue,
urlcolor=MidnightBlue,pdftitle={Omnific Points of Algebraic Groups},
pdfauthor={OpenAI ChatGPT},pdfsubject={Unipotent rigidity, Gaussian omnific points, and universal elementary quotients}}
\usepackage[nameinlink,noabbrev]{cleveref}
\setlength{\parskip}{3pt}
\setlength{\parindent}{1.15em}
\setlength{\emergencystretch}{2em}
\newcolumntype{P}[1]{>{\raggedright\arraybackslash}p{#1}}
\setlist[enumerate]{leftmargin=2em,itemsep=3pt,topsep=4pt}
\pagestyle{fancy}\fancyhf{}
\fancyhead[L]{\small Omnific points and elementary quotients}
\fancyhead[R]{\small Research article}
\fancyfoot[C]{\thepage}
\renewcommand{\headrulewidth}{0.3pt}
\numberwithin{equation}{section}
\newtheorem{theorem}{Theorem}[section]
\newaliascnt{lemma}{theorem}
\newtheorem{lemma}[lemma]{Lemma}
\aliascntresetthe{lemma}
\newaliascnt{proposition}{theorem}
\newtheorem{proposition}[proposition]{Proposition}
\aliascntresetthe{proposition}
\newaliascnt{corollary}{theorem}
\newtheorem{corollary}[corollary]{Corollary}
\aliascntresetthe{corollary}
\newtheorem{mainthm}{Theorem}
\renewcommand{\themainthm}{\Alph{mainthm}}
\theoremstyle{definition}
\newaliascnt{definition}{theorem}
\newtheorem{definition}[definition]{Definition}
\aliascntresetthe{definition}
\newaliascnt{example}{theorem}
\newtheorem{example}[example]{Example}
\aliascntresetthe{example}
\theoremstyle{remark}
\newaliascnt{remark}{theorem}
\newtheorem{remark}[remark]{Remark}
\aliascntresetthe{remark}
\newaliascnt{question}{theorem}
\newtheorem{question}[question]{Question}
\aliascntresetthe{question}
\crefname{mainthm}{Theorem}{Theorems}
\crefname{theorem}{Theorem}{Theorems}
\crefname{lemma}{Lemma}{Lemmas}
\crefname{proposition}{Proposition}{Propositions}
\crefname{corollary}{Corollary}{Corollaries}
\newcommand{\No}{\mathbf{No}}
\newcommand{\Oz}{\mathbf{Oz}}
\newcommand{\GOz}{\mathbf{Oz}[i]}
\newcommand{\R}{\mathbb R}\newcommand{\C}{\mathbb C}
\newcommand{\Z}{\mathbb Z}\newcommand{\Q}{\mathbb Q}\newcommand{\N}{\mathbb N}
\newcommand{\I}{\mathcal I}
\newcommand{\Bk}{\mathcal B_k}\newcommand{\BR}{\mathcal B_{\R}}
\newcommand{\BC}{\mathcal B_{\C}}
\newcommand{\cG}{\mathcal G}
\newcommand{\Ga}{\mathbb G_a}\newcommand{\Gm}{\mathbb G_m}
\newcommand{\ct}{\operatorname{ct}}
\newcommand{\supp}{\operatorname{supp}}
\newcommand{\GL}{\operatorname{GL}}\newcommand{\SL}{\operatorname{SL}}
\newcommand{\SO}{\operatorname{SO}}\newcommand{\OO}{\operatorname{O}}
\newcommand{\UT}{\operatorname{UT}}
\newcommand{\Lie}{\operatorname{Lie}}\newcommand{\Hom}{\operatorname{Hom}}
\newcommand{\Spec}{\operatorname{Spec}}\newcommand{\tr}{\operatorname{tr}}
\newcommand{\Id}{\operatorname{id}}\newcommand{\Sym}{\operatorname{Sym}}
\newcommand{\der}{\mathrm{der}}
\newcommand{\ee}[2]{e_{#1}(#2)}
\newcommand{\bdeg}{\operatorname{deg}_{\omega}}
\newcommand{\pin}{\texttt{220784a4766b9dd9ab3deb92f9f28ddfdd866187}}
\newcommand{\status}[1]{\par\smallskip\noindent\textit{Status and provenance.} #1\par\smallskip}

\begin{document}
\begin{titlepage}
\centering
\vspace*{4mm}
{\large SURREAL AND SURCOMPLEX ALGEBRA\par}
\vspace{7mm}
{\huge\bfseries Omnific Points\\[3mm]of Algebraic Groups\par}
\vspace{6mm}
{\Large A unipotent dichotomy and\\[2mm]universal elementary quotients\par}
\vspace{8mm}
{\large Prepared for the Surreal project\par}
\vspace{4mm}
{23 September 2026\par}
\vspace{7mm}
\begin{minipage}{0.92\textwidth}
\begin{abstract}
We study algebraic groups over the omnific integers and over their Gaussian
complexification. A real algebraic group has no nonconstant points over the
ring of surreal normal forms with nonnegative exponents precisely when it
has no nonidentity real unipotent element. Equivalently, its identity
component is reductive with compact real derived group. For every fixed
integral model this gives an exact dichotomy: either all omnific points
are ordinary integral points, or the constant-term kernel contains a copy
of the entire additive class of purely infinite normal forms. Over the
Gaussian omnific integers the criterion changes: the complex identity
component must be a torus.

The second theorem concerns arbitrary abstract representations, not just
algebraic or continuous ones. For every ordinary integer $n\ge3$, the subgroup
of $E_n(\Oz)$ generated by purely infinite elementary matrices is perfect,
normal, and has no nontrivial homomorphism into any set-sized group. Thus
constant term is the universal set-sized group quotient of $E_n(\Oz)$,
with target $\SL_n(\Z)$. The analogous target over $\Oz[i]$ is
$\SL_n(\Z[i])$. These kernels are not simple: an explicit proper-class
chain of congruence subgroups separates their elements, although no
set-sized subfamily does.

The proofs combine normal-form support arguments, classical algebraic-group
structure, and finite matrix commutators. Explicit orthogonal examples,
unipotent kernel coordinates, unitriangular filtrations, and fixed-workspace
limitations are included. The new theorem package is separated from the
repository's earlier norm-rigidity and ring-quotient results.
\end{abstract}
\end{minipage}
\vfill
\begin{minipage}{0.92\textwidth}\small
\textbf{Research status.} This is an AI-assisted, proof-bearing research draft.
The main combined statements are proposed contributions, not claims of
verified historical priority. They were not located in the inspected
sources. Classical inputs are attributed; all new deductions are proved.
The accompanying symbolic checks are not a Lean formalization or a substitute
for independent mathematical review. No named longstanding conjecture is
claimed to be settled.
\end{minipage}
\end{titlepage}
\begingroup
\small
\setlength{\parskip}{0pt}
\tableofcontents
\endgroup
\newpage

\section{Introduction and the theorem package}\label{sec:intro}

The omnific integers are not a nonstandard model of ordinary integer
arithmetic. Their positive-exponent coefficients are arbitrary real numbers,
while only the constant coefficient is required to be an ordinary integer.
This mixture makes algebraic geometry over them quite different from either
geometry over $\Z$ or geometry over the full surreal field.

The repository already establishes the utility of the retraction
$\ct:\Oz\to\Z$, develops norm and quadratic rigidity, and proves a
universal set-sized \emph{ring} quotient theorem
\cite{RepoDio,RepoQuot}. We ask two further questions. Which algebraic groups
acquire new omnific points? And can the infinite directions of an elementary
matrix group act on any ordinary set, even when no algebraicity is imposed
on that action? The answers reveal two distinct mechanisms: geometric
rigidity controlled by unipotents, and an abstract representation obstruction
controlled by the supply of surreal exponent scales.

\subsection{Notation and an elementary warning}
For $k=\R$ or $\C$, let
\[
 \I_k=\left\{\sum_{s\in S}a_s\omega^s:
       S\subset\No^{>0}\text{ is a reverse-well-ordered set},\ a_s\in k\right\}.
\]
Zero is included. Put
\begin{equation}\label{eq:mainrings}
 \Bk=k\oplus\I_k,\qquad A_{D,k}=D\oplus\I_k\quad(D\subseteq k).
\end{equation}
Then $A_{\Z,\R}=\Oz$ and $A_{\Z[i],\C}=\GOz$.
These are additive direct sums, not ring direct products. The map $\ct$
extracts the coefficient of $\omega^0$.

The ideal $\I_k$ is \emph{purely infinite}, not infinitesimal. For example,
$\sqrt2\,\omega$ and $-\omega^\omega$ belong to $\I_{\R}$, whereas
$\omega^{-1}$ does not. Coefficient extraction on these rings must not be
confused with standard part on the valuation ring of finite surreals.
All group ranks, matrix sizes, and polynomial degrees below are ordinary
finite integers.

\subsection{The geometric dichotomy}
An algebraic group is affine and of finite type unless stated otherwise.
For a real algebraic group $G$, write $G^0$ for its algebraic identity
component, not its identity component as a real Lie group.

\begin{mainthm}[Omnific rigidity and real unipotents]\label{main:real}
For a real algebraic group $G$, the following are equivalent:
\begin{enumerate}
\item $G(\BR)=G(\R)$.
\item $G(\R)$ has no unipotent element other than the identity.
\item $G$ has no real algebraic subgroup isomorphic to $\Ga$.
\item $G^0$ is reductive and $(G^0)^{\der}(\R)$ is compact.
\end{enumerate}
Let $\cG\hookrightarrow\GL_{N,\Z}$ be a closed subgroup scheme and
$G=\cG_{\R}$. These conditions are also equivalent to
\[
                         \cG(\Oz)=\cG(\Z).
\]
When they fail, there is an injective group homomorphism
\[
 (\I_{\R},+)\longrightarrow
 \ker\bigl(\ct:\cG(\Oz)\longrightarrow\cG(\Z)\bigr),
 \qquad a\longmapsto\exp(aN),
\]
for a fixed nonzero real nilpotent matrix $N$. The exponential is a finite
polynomial. Every ordinary integral point lies on a corresponding
injectively parameterized family of omnific points.
\end{mainthm}

Compactness of the whole real group is not the criterion. Every real torus,
including $\Gm$, is rigid. A noncompact semisimple real group, in contrast,
is not. Nor can real unipotents be replaced by rational unipotents:
\cref{sec:examples} gives an integral orthogonal group with no nonidentity
rational unipotents but an explicit purely infinite unipotent family.

\begin{mainthm}[The Gaussian dichotomy]\label{main:complex}
For a complex algebraic group $G$,
\[
 G(\BC)=G(\C)
 \quad\Longleftrightarrow\quad
 G^0\text{ is a torus}.
\]
For a closed subgroup scheme $\cG\hookrightarrow\GL_{N,\Z[i]}$,
\[
 \cG(\GOz)=\cG(\Z[i])
 \quad\Longleftrightarrow\quad
 (\cG_{\C})^0\text{ is a torus}.
\]
If rigidity fails, the constant-term kernel contains an injective image
of $(\I_{\C},+)$ under a finite matrix exponential.
\end{mainthm}

For example, $\SO_3(\Oz)$ consists of the $24$ integral rotations, but
$\SO_3(\GOz)$ contains nonconstant one-parameter unipotent families.
The group scheme $\SO_3$ is the same; the change of coefficient field
changes the answer.

\subsection{The abstract group quotient}
For any commutative unital ring $A$, $E_n(A)$ denotes the subgroup of
$\SL_n(A)$ generated by the elementary matrices
\[
 e_{ij}(a)=1_n+aE_{ij},\qquad a\in A,\quad i\ne j.
\]
Generation always means finite words. We do not identify $E_n(A)$ with
$\SL_n(A)$ for an omnific ring.

\begin{mainthm}[Universal set-sized elementary quotient]\label{main:elementary}
Let $(D,k)=(\Z,\R)$ or $(\Z[i],\C)$, put $A=A_{D,k}$, and let $n\ge3$.
Define
\[
 P_n(k)=\langle e_{ij}(a):a\in\I_k,\ i\ne j\rangle.
\]
Then $P_n(k)$ is a nontrivial perfect normal subgroup of $E_n(A)$, and
\begin{equation}\label{eq:introsplit}
 E_n(A)=P_n(k)\rtimes\SL_n(D),\qquad
 P_n(k)=\ker\bigl(\ct:E_n(A)\to\SL_n(D)\bigr).
\end{equation}
Every class group homomorphism from $P_n(k)$ to a set-sized group is trivial.
Consequently, for every set-sized group $H$, composition with $\ct$ gives
an exact correspondence
\[
 \Hom\bigl(\SL_n(D),H\bigr)
 \ \longleftrightarrow\ 
 \Hom\bigl(E_n(A),H\bigr).
\]
No continuity, preservation of infinite sums, definability in the ring
language, or algebraicity of a homomorphism is assumed.
\end{mainthm}

The assertion about $P_n(k)$ is stronger than a factorization theorem for
homomorphisms defined on $E_n(A)$: a homomorphism defined only on the
kernel need not extend to the whole group. Our proof treats it directly.
The ring-quotient theorem in \cite{RepoQuot} motivates the result and supplies
the geometric-series collision mechanism. Passing to arbitrary abstract
group maps requires the root commutator argument developed here.

\subsection{What is proposed as new}
The intended contribution is the complete group-level dichotomy of
\cref{main:real,main:complex}, together with the abstract kernel and universal
quotient in \cref{main:elementary} and the congruence tower in
\cref{sec:tower}. These combined statements were not found in the sources
inspected for this article. This is a qualified novelty assessment, not a
proof of priority.

Normal forms, real closedness, torus structure, nilpotent exponentials,
Cartan theory, and elementary commutator identities are classical. The
constant-term formalism, common monomial divisibility, and geometric collision
identities already occur in the repository. We reprove the needed forms
rather than import their correctness from an unreviewed draft. The
unitriangular and Lie-algebra descriptions are useful specializations of
standard machinery; they are not presented as independent major discoveries.

\section{Normal forms, constant terms, and class size}\label{sec:rings}

\subsection{The normal-form input}
The standard normal-form theorem identifies $\No$ with real Hahn normal
forms
\begin{equation}\label{eq:nf}
                    f=\sum_{s\in S}a_s\omega^s,
\end{equation}
where $S\subset\No$ is a set and is reverse well ordered: every nonempty
subset of $S$ has a greatest element. The omega-map satisfies
$\omega^{s+t}=\omega^s\omega^t$. Addition is coefficientwise and
multiplication is Hahn convolution. The sign of a nonzero real normal form
is the sign of its leading coefficient. These facts, and real closedness
of $\No$, are standard \cite{Conway,Gonshor,Berarducci}.

Complex coefficients give exactly $\No[i]$, by separating real and imaginary
parts. A finite union of reverse-well-ordered supports is reverse well ordered,
so this complexification needs no larger support convention. The notation
$\omega^s$ is the Conway monomial, not an unspecified analytic exponential.
In the common convention $t=\omega^{-1}$, all exponent signs are reversed.

We use the usual Hahn support lemma: products of Hahn series are defined,
and a series supported on strictly negative exponents has a well-defined
formal geometric inverse near $1$. Thus, for $\gamma>0$,
\begin{equation}\label{eq:geom}
 (1-\omega^{-\gamma})^{-1}=\sum_{m\ge0}\omega^{-m\gamma}.
\end{equation}
This is a formal identity of normal forms. It does not assert convergence
of a sequence in the fine topology of the full surreal field.

\subsection{The elementary ring facts}
\begin{lemma}[Retraction and units]\label{lem:units}
For $k=\R,\C$ and a unital subring $D\subseteq k$, $\I_k$ is an ideal
of $\Bk$ and of $A_{D,k}$. The map $\ct:A_{D,k}\to D$ is a split
surjective ring homomorphism with kernel $\I_k$. Moreover,
\[
                       \Bk^\times=k^\times,
 \qquad A_{D,k}^\times=D^\times.
\]
\end{lemma}
\begin{proof}
All exponents in these rings are nonnegative. An exponent sum is zero
only when both summands are zero. This proves multiplicativity of
constant coefficient and the ideal assertion.

For nonzero $f\in\Bk$, let $\bdeg(f)=\max\supp(f)$. Hahn multiplication
has nonzero leading coefficient equal to the product of the leading
coefficients, so
\[
 \bdeg(fg)=\bdeg(f)+\bdeg(g),\qquad\bdeg(f)\ge0.
\]
If $fg=1$, both degrees are zero. A normal form with nonnegative support
and leading exponent zero is constant. Hence every unit in $\Bk$ is a
nonzero constant. Intersecting with $A_{D,k}$ and applying the same argument
to the inverse gives the last assertion.
\end{proof}

\begin{lemma}[Bounded coordinates are constant]\label{lem:bounded}
If $f\in\BR$ and $|f|\le M$ for an ordinary real $M$, then $f\in\R$.
Consequently $\Oz\cap\R=\Z$. If $z\in\BC$ has real and imaginary
parts bounded by ordinary real numbers, then $z\in\C$.
\end{lemma}
\begin{proof}
If $f$ has a positive leading exponent $s$, its absolute value exceeds
every ordinary real: the positive-exponent leading term dominates any
constant. This contradicts the bound. The remaining assertions follow
coordinatewise and from the definition of the constant coefficient.
\end{proof}

For example, $\omega+\omega^{-1}$ does not belong to $\BR$, despite being
infinite. It is the absence of negative exponents, not merely the size of
an element, that makes \cref{lem:units,lem:bounded} work.

\subsection{Set-theoretic conventions}
We work in G\"odel--Bernays set theory with global choice, or an equivalent
standard set/class framework for surreal numbers. Each individual surreal,
normal-form support, matrix, and finite word is a set. The full rings and
several groups considered here are proper classes. A homomorphism involving
a class group means a class function satisfying the usual finite group
identities. We never require a set of all such maps.

Statements about a class quotient may be read as statements about a kernel
and a universal property. If a literal quotient carrier is desired, Scott's
trick gives set-valued codes for its equivalence classes. No argument uses
a cardinal number assigned to a proper class. Instead, a target set $H$
is contradicted by constructing a set of more than $|H|$ source elements.

\begin{lemma}[A gap below every set of positive exponents]\label{lem:gap}
For every set $S\subset\No^{>0}$ there is $\delta>0$ with
$\delta<s$ for all $s\in S$. Every set of surreal numbers has a strict
upper bound.
\end{lemma}
\begin{proof}
For nonempty $S$, the separated cut $\{0\mid S\}$ gives the desired
number. Use $\delta=1$ when $S$ is empty. A strict upper bound for a set
$T$ is supplied by $\{T\mid\varnothing\}$.
\end{proof}

\begin{proposition}[Single-product divisibility]\label{prop:factor}
Every $c\in\I_k$ is a product of two members of $\I_k$. More generally,
for every set-sized family $\mathcal F\subseteq\I_k$ there is $\delta>0$
such that $\omega^{-\delta}f\in\I_k$ for every $f\in\mathcal F$.
In particular, $\I_k^r=\I_k$ for every ordinary $r\ge1$.
\end{proposition}
\begin{proof}
The union of the supports in $\mathcal F$ is a set. Choose $\delta$
strictly below its positive members by \cref{lem:gap}. Translating supports
by $-\delta$ leaves them strictly positive. For an individual element,
\[
                   c=\omega^\delta\bigl(\omega^{-\delta}c\bigr).
\]
Both factors lie in $\I_k$. Zero causes no exception. Induction gives
the assertion about powers.
\end{proof}

\status{The retraction, unit argument, and common-support divisibility are
background or already present in \cite{RepoDio,RepoQuot}. The full-class
hypothesis in \cref{lem:gap} is substantive. For example, the positive
support $\{1/m:m\ge1\}$ has no positive lower bound in the exponent
group $\R$, but has one in $\No$.}

\section{A bounded-fiber rigidity principle}\label{sec:fibers}

\subsection{Tori and finite schemes}
A torus over $k$ is an algebraic group that becomes a product of copies
of $\Gm$ over an algebraic closure of $k$.

\begin{lemma}[Torus freezing]\label{lem:torus}
For every real torus $T$, $T(\BR)=T(\R)$. For every complex torus $T$,
$T(\BC)=T(\C)$.
\end{lemma}
\begin{proof}
A complex torus is isomorphic over $\C$ to $\Gm^r$. Its points over
$\BC$ are $(\BC^\times)^r=(\C^\times)^r$ by \cref{lem:units}.

For a real torus, extend scalars from $\BR$ to
$\BR\otimes_{\R}\C=\BC$. After a complex splitting, the resulting
point is constant by the preceding paragraph. Hence the original point
is a constant complex point fixed by the real conjugation descent datum.
Such points are exactly $T(\R)$. This is ordinary Galois descent for
points; the inclusion $\BR\hookrightarrow\BC$ is injective.
\end{proof}

\begin{lemma}[Finite schemes do not acquire points]\label{lem:finite}
If $X$ is a finite scheme over $k=\R$ or $\C$, then $X(\Bk)=X(k)$.
\end{lemma}
\begin{proof}
Every element of the finite-dimensional coordinate algebra $k[X]$ is
algebraic over $k$. Its image under a point with values in $\Bk$ is
therefore algebraic over $k$ in the ambient field. The only elements of
$\No$ algebraic over $\R$ are real, because $\R$ is real closed and
$\No$ is ordered. The corresponding complex assertion follows because
$\C$ is algebraically closed. Thus every coordinate image is constant.
Nilpotents in $k[X]$ map to zero and do not affect the argument.
\end{proof}

The same proof applies to any zero-dimensional affine scheme of finite
type over $k$: its coordinate algebra is Artinian and finite-dimensional.
We use this observation later without a reducedness assumption.

\subsection{Bounded real fibers}
\begin{theorem}[Bounded fibers above a torus]\label{thm:fibers}
Let $X\hookrightarrow\mathbb A_{\R}^m$ be an affine real scheme of
finite type, let $T$ be a real torus, and let $\pi:X\to T$ be a
morphism. Suppose that for every $t\in T(\R)$ the real points of the
fiber $X_t$ are bounded in the specified affine coordinates. Then
\[
                              X(\BR)=X(\R).
\]
No bound uniform in $t$ is required.
\end{theorem}
\begin{proof}
Take $x\in X(\BR)$. By \cref{lem:torus}, $\pi(x)$ is a constant point
$t\in T(\R)$. Applying $\ct$ gives a real point $\ct(x)$ of $X_t$,
so this fiber is not empty.

Choose an ordinary real number $M>0$ bounding all real coordinates of
$X_t(\R)$. The statement that every solution of the finite polynomial
equations for $X_t$ satisfies $-M\le x_j\le M$ is a first-order formula
in the language of ordered fields with real parameters. Every embedding
of real closed fields is elementary. Apply this standard consequence of
quantifier elimination to $\R\subset K$, where $K\subset\No$ is a
set-sized real closed subfield containing $\R$ and all coordinates of
$x$. Such a $K$ is obtained by taking a real closure inside $\No$ of
the field generated by this set. The bounds therefore hold for $x$.
By \cref{lem:bounded}, each coordinate of $x$ is real.
\end{proof}

The only transfer used here is transfer between real closed \emph{fields}.
We do not assert elementary equivalence of $\Z$ and $\Oz$, nor transfer
in a language containing an omnific-integer predicate. The model-theoretic
input is standard; see, for example, \cite{Marker}.

\begin{corollary}[Complex zero-dimensional fibers]\label{cor:quasifinite}
If $X$ is affine of finite type over $\C$ and $\pi:X\to T$ is a
morphism to a complex torus whose fibers over complex points are
zero-dimensional, then $X(\BC)=X(\C)$.
\end{corollary}
\begin{proof}
Freeze $\pi(x)$ by \cref{lem:torus}, then apply the zero-dimensional
version of \cref{lem:finite} to that fiber.
\end{proof}

\begin{corollary}[Compact-by-torus groups]\label{cor:compacttorus}
Let $G$ be a connected real reductive algebraic group. If
$H=G^{\der}$ has compact real points, then $G(\BR)=G(\R)$.
\end{corollary}
\begin{proof}
The quotient $T=G/H$ is a real torus, a standard structure theorem for
connected reductive groups \cite{MilneAG}. A fiber with a real point $g$
has real points exactly $gH(\R)$, hence is compact. In any affine
embedding its coordinate functions are bounded. Apply
\cref{thm:fibers}. Empty real fibers also satisfy the boundedness condition.
\end{proof}

\status{The bounded-fiber formulation is the bridge developed here between
normal-form rigidity and algebraic groups. Its inputs are unit rigidity,
classical torus splitting, and real-closed-field transfer. It generalizes the
kind of bounded-coordinate and constant-product arguments in the repository,
without treating those earlier arguments as new.}

\section{Unipotents and the complete group classification}\label{sec:dichotomy}

\subsection{A polynomial subgroup from a single unipotent}
A matrix $u\in\GL_N(k)$ is unipotent if $u-1_N$ is nilpotent. This
notion is independent of the faithful representation of an algebraic
group. The following construction will also control integrality.

\begin{lemma}[Polynomial one-parameter subgroup]\label{lem:exp}
Let $k$ be a characteristic-zero field, let $G\subseteq\GL_{N,k}$ be
an algebraic group, and let $1_N\ne u\in G(k)$ be unipotent. There is a
nonzero nilpotent $N_0\in M_N(k)$ such that
\[
 N_0=\log u=\sum_{j=1}^{N-1}\frac{(-1)^{j+1}}j(u-1_N)^j,
 \qquad v(t)=\exp(tN_0)=\sum_{j=0}^{N-1}\frac{t^jN_0^j}{j!}
\]
defines a closed subgroup $v:\Ga\hookrightarrow G$.
For any $k$-algebra, $v(s+t)=v(s)v(t)$ and $v(-t)=v(t)^{-1}$.
\end{lemma}
\begin{proof}
All exponential and logarithm identities here are identities modulo a
fixed nilpotence power; thus they are polynomial identities over $\Q$.
In particular, $N_0$ is nilpotent, $\exp N_0=u$, and $N_0\ne0$.
For each nonnegative integer $m$, $v(m)=u^m\in G(k)$.
The matrix $v(t)$ has determinant $1$. Every defining equation of $G$
in the affine coordinates of $\GL_N$ therefore pulls back to a polynomial
in $k[t]$ vanishing at all nonnegative integers. It is zero. Hence the
whole polynomial map lands in $G$.

Choose an entry $(p,q)$ with $(N_0)_{pq}\ne0$. On the image, the
polynomial function
\[
             g\longmapsto\frac{(\log g)_{pq}}{(N_0)_{pq}}
\]
recovers $t$, where $\log g$ is the same fixed truncated polynomial in
$g-1_N$. The resulting map $k[G]\to k[t]$ is surjective. This proves
that $v$ is a closed immersion, as well as injectivity on points. The
addition and inverse formulas follow from the commuting nilpotent
exponential identities.
\end{proof}

\begin{corollary}[Purely infinite parameters]\label{cor:expI}
Let $k=\R$ or $\C$, and let $D=\Z$ or $\Z[i]$ respectively. If the
$k$-fiber of a closed subgroup scheme $\cG\subseteq\GL_{N,D}$ contains
a nonidentity unipotent, then the map of \cref{lem:exp} restricts to
\[
 v:(\I_k,+)\hookrightarrow\cG(A_{D,k}),\qquad \ct(v(a))=1_N.
\]
\end{corollary}
\begin{proof}
For $a\in\I_k$, every positive-degree term in $v(a)$ has entries in
$\I_k$, since $\I_k$ is closed under multiplication and multiplication
by $k$. Thus $v(a)$ has entries in $D+\I_k$, and so does $v(-a)$.
All the original $D$-equations vanish: their substitution into $v(t)$
is the zero polynomial over $k$, as in \cref{lem:exp}. The constant term
is the identity. Injectivity follows either from the displayed inverse
coordinate or from $\log(v(a))=aN_0$.
\end{proof}

This explains why rationality of the subgroup is unnecessary. Irrational
coefficients in $N_0$ multiply purely infinite terms and do not alter the
integral constant coefficient.

\subsection{The real structural input, with a short reduction}
We record precisely the classical group structure needed. A connected
linear algebraic group in characteristic zero has a unipotent radical
$R_u(G)$; it is reductive when this radical is trivial. For a connected
reductive group, the derived subgroup is connected semisimple and the
quotient by it is a torus. A nontrivial unipotent group over $\R$ has
nontrivial real points; its exponential is a polynomial isomorphism from
its nilpotent Lie algebra as a variety. These are standard facts, including
\cite[Proposition 14.32]{MilneAG} for the last statement.

For completeness, the compactness step can be reduced to elementary
matrix arguments once Cartan's theorem is imported.

\begin{lemma}[Real semisimple compactness criterion]\label{lem:compact}
A connected semisimple real algebraic group $H$ has no nonidentity real
unipotent element if and only if $H(\R)$ is compact.
\end{lemma}
\begin{proof}
If $H(\R)$ is compact, average an inner product in a faithful real
representation. In the resulting basis every element is orthogonal.
An orthogonal unipotent matrix is the identity, because orthogonal
matrices are diagonalizable over $\C$.

Conversely, Cartan's theorem permits a basis of a faithful representation
in which $H$ is stable under transpose. This is the standard matrix
form of the theorem; see \cite[Theorem 1.16 and Example 1.17(c)]{MilneSV}.
Its real Lie algebra $\mathfrak h$ is therefore stable under transpose.
Suppose first that every element of $\mathfrak h$ is skew-symmetric.
The real Lie identity component of $H(\R)$ lies in the orthogonal group
and is closed there, hence compact. A real algebraic group has finitely
many real connected components, so $H(\R)$ is compact.

If $H(\R)$ is not compact, there is consequently a nonzero real symmetric
matrix $S\in\mathfrak h$. The Lie algebra is semisimple, so its center
is zero and $\operatorname{ad}S$ is nonzero. On $\mathfrak h$, with inner
product $\langle A,B\rangle=\tr(AB^T)$, the operator
$A\mapsto[S,A]$ is self-adjoint: the equality follows by cyclically
permuting factors inside the trace. It has an eigenvector $Y\ne0$ with
a nonzero real eigenvalue $\lambda$.

Diagonalize $S$ orthogonally, with eigenvalues $s_1,\ldots,s_N$. The
identity $[S,Y]=\lambda Y$ implies that $Y$ sends the $s$-eigenspace of
$S$ to its $(s+\lambda)$-eigenspace. Repeated applications must leave
the finite eigenvalue set, so $Y$ is nilpotent. The usual exponential
map of the real matrix Lie group sends $tY$ into $H(\R)$. Here that
exponential is a finite polynomial. Its value at $t=1$ is a nonidentity
unipotent, since the truncated logarithm recovers $Y$. This contradicts
the assumed absence of unipotents.
\end{proof}

The proof uses ordinary real Lie theory only at the Cartan and Lie
exponential steps. It does not invoke a global exponential on the
surcomplex field. Equivalently one can obtain the noncompact direction
from real rank and relative root groups \cite[Sections 22--23]{MilneRG}.

\subsection{Proof of the real dichotomy}
\begin{proof}[Proof of \cref{main:real}]
Assume first that $G(\R)$ contains no nonidentity unipotent. The
unipotent radical of $G^0$ is then trivial, so $G^0$ is reductive.
Its derived group $H$ also has no real unipotents, hence has compact real
points by \cref{lem:compact}. By \cref{cor:compacttorus},
$G^0(\BR)=G^0(\R)$.

Take any $g\in G(\BR)$ and let $g_0=\ct(g)\in G(\R)$. The component
group scheme $\pi_0(G)=G/G^0$ is finite in characteristic zero. Its
$\BR$-points are constant by \cref{lem:finite}. Thus $g$ and $g_0$
have the same image there, and $g_0^{-1}g\in G^0(\BR)$. This element
is real; applying $\ct$ shows that it equals the identity. Hence
$g=g_0\in G(\R)$.

If there is a nonidentity unipotent, \cref{lem:exp} and
\cref{cor:expI} (or the same argument over $\BR$) supply nonconstant
points, so rigidity fails. The same construction gives a real $\Ga$
subgroup. Conversely a $\Ga$ subgroup gives nonconstant points by
substituting a purely infinite parameter into its nonconstant polynomial
coordinates. Thus conditions (1), (2), and (3) are equivalent.

We have already proved that (2) implies (4), and (4) implies (1) by
\cref{cor:compacttorus} and the component argument. This completes the
four-way equivalence.

For an integral model, the inclusion into its real fiber and the first
part show that an omnific point is constant in the rigid case. Its matrix
entries and determinant-inverse coordinate then lie in $\Z$, by
\cref{lem:units,lem:bounded}; it is an ordinary integral point of the
same model. Conversely every integral point is an omnific point. In the
nonrigid case \cref{cor:expI} gives the claimed injection. Left translation
by any $q\in\cG(\Z)$ gives the family $qv(a)$ through $q$ and preserves
injectivity.
\end{proof}

Only the characteristic-zero fiber is used. We do not need smoothness
of the integral model at every prime. Keeping the model fixed is important:
ordinary integral points do depend on that model, whereas the criterion
for acquiring nonconstant points depends only on the real fiber.

\subsection{Proof of the Gaussian dichotomy}
\begin{proof}[Proof of \cref{main:complex}]
Over an algebraically closed field of characteristic zero, a connected
affine algebraic group with no nontrivial connected unipotent subgroup
is a torus \cite[Corollary 17.25]{MilneAG}. Equivalently, a connected
group that is not a torus has a nonidentity unipotent. For example this
follows by first removing the unipotent radical and then using a root
subgroup of a nontrivial semisimple derived group.

If $G^0$ is a torus, \cref{lem:torus} freezes its points. The finite
component argument in the real proof gives $G(\BC)=G(\C)$. Otherwise
\cref{lem:exp} supplies an injective purely infinite family. For an
integral Gaussian model, constant points in $\GOz$ have entries in
$\Z[i]$, and \cref{cor:expI} supplies nonconstant integral-model points
in the other case. This proves all assertions.
\end{proof}

\section{Explicit orthogonal examples}\label{sec:examples}

\subsection{Compact real groups and Gaussian escape}
For the standard positive definite form, an orthogonal matrix $g$ over
$\Oz$ satisfies $\sum_i g_{ij}^2=1$ for each column. Each entry lies
between $-1$ and $1$ and is therefore an ordinary integer. Its columns
are signed coordinate vectors. Consequently
\begin{equation}\label{eq:orthogonal}
 \OO_n(\Oz)=\OO_n(\Z)
   =\{\text{signed permutation matrices}\}.
\end{equation}
There are $2^n n!$ such matrices and, for $n\ge1$, exactly half have
determinant $1$. In particular $|\SO_3(\Oz)|=24$.

Over $\C$, take
\[
 v=(1,i,0)^T,\qquad w=(0,0,1)^T,
 \qquad N=vw^T-wv^T
       =\begin{pmatrix}0&0&1\\0&0&i\\-1&-i&0\end{pmatrix}.
\]
Then $N^T=-N$, $N^3=0$, and $N^2\ne0$. The finite exponential is
\begin{equation}\label{eq:complexexample}
 g(a)=1_3+aN+\frac{a^2}{2}N^2
 =\begin{pmatrix}
 1-a^2/2&-ia^2/2&a\\
 -ia^2/2&1+a^2/2&ia\\
 -a&-ia&1
 \end{pmatrix}.
\end{equation}
For every $a\in\I_{\C}$, this is in $\SO_3(\GOz)$, has constant
term $1_3$, and is nonconstant when $a\ne0$. Orthogonality follows from
$\exp(aN)^T=\exp(-aN)$; the determinant is $1$ because $N$ is
nilpotent. The $(1,3)$ entry recovers $a$.

Thus neither compactness of the real form nor a positive-definite defining
form prevents new \emph{Gaussian} omnific points. The equation
$z_1^2+z_2^2+z_3^2=1$ is not a positivity statement over $\C$.
In contrast, a conjugate-transpose condition defines a real algebraic
unitary group and must be analyzed as a real group, not as this complex
orthogonal group. This distinction prevents a common misreading of the
example.

\subsection{A rationally anisotropic group with omnific unipotents}
Consider
\[
 Q=\operatorname{diag}(1,1,-3),\qquad q(x,y,z)=x^2+y^2-3z^2.
\]
The form is anisotropic over $\Q$. Indeed, after clearing denominators,
a primitive integer solution to $x^2+y^2=3z^2$ would have $x$ and $y$
divisible by $3$, since a square modulo $3$ is $0$ or $1$. The equation
then forces $3\mid z$, contradicting primitivity.

There is no nonidentity unipotent in $\SO(q)(\Q)$. To see this without
using a rank classification, suppose $u$ were such a matrix and put
$N=\log u\ne0$. By \cref{lem:exp}, $\exp(tN)$ preserves $q$ for
all $t$, so differentiating this polynomial identity at zero gives
$N^TQ+QN=0$. Choose $r\ge1$ maximal with $N^r\ne0$, and a rational
vector $z$ with $N^rz\ne0$. For the bilinear form $B(x,y)=x^TQy$,
\[
 B(N^rz,N^rz)=(-1)^rB(z,N^{2r}z)=0.
\]
This contradicts rational anisotropy.

Nevertheless, put $v=(\sqrt3,0,1)^T$ and $w=(0,1,0)^T$. Then
$B(v,v)=B(v,w)=0$, $B(w,w)=1$, and
\[
 N=vw^TQ-wv^TQ
  =\begin{pmatrix}0&\sqrt3&0\\-\sqrt3&0&3\\0&1&0\end{pmatrix}
\]
satisfies $N^TQ+QN=0$ and $N^3=0$. Explicitly,
\begin{equation}\label{eq:realexample}
 \exp(aN)=
 \begin{pmatrix}
 1-3a^2/2&\sqrt3a&3\sqrt3a^2/2\\
 -\sqrt3a&1&3a\\
 -\sqrt3a^2/2&a&1+3a^2/2
 \end{pmatrix}.
\end{equation}
For each $a\in\I_{\R}$ this is in $\SO(q)(\Oz)$ and has constant
term $1_3$. Its $(3,2)$ entry is $a$, so the parameterization is injective.
In particular, the single substitution $a=\omega$ produces a fully explicit
nonconstant omnific matrix.

This example isolates the role of the coefficient field: the polynomial
subgroup exists over $\R$, not over $\Q$, and that is sufficient for
omnific integrality. It does not contradict the repository's rigidity
of nonzero binary norm fibers; here the ambient group is three-dimensional
and its real derived group is noncompact.

\subsection{A comparison table}
\begin{center}
\begin{tabular}{P{0.24\textwidth}P{0.31\textwidth}P{0.34\textwidth}}
\toprule
Group or model & Real omnific behavior & Reason \\
\midrule
A real torus & All points constant & Units freeze after complex splitting \\
Compact semisimple group & All points constant & Bounded real coordinates \\
$\SO(x^2-y^2)$ & All points constant & A torus, although noncompact \\
$\SL_n$, $n\ge2$ & Purely infinite families & Elementary $\Ga$ subgroups \\
$\SO(x^2+y^2-3z^2)$ & Purely infinite families & Real, not rational, unipotents \\
\bottomrule
\end{tabular}
\end{center}
The entries concern constant-term rigidity, not finite generation of an
arithmetic subgroup or elementary generation of all omnific matrices.

\section{Unipotent kernels and exact unitriangular filtrations}\label{sec:unipotent}

\subsection{Finite logarithmic coordinates}
Let $U$ be a unipotent algebraic group over $k=\R$ or $\C$. In a faithful
representation it is conjugate over $k$ to a group of upper unitriangular
matrices. Its Lie algebra $\mathfrak u$ consists of nilpotent matrices
with a common finite nilpotence bound. The usual exponential and logarithm
are mutually inverse polynomial maps; this is the classical
characteristic-zero unipotent correspondence \cite[Proposition 14.32]{MilneAG}.

\begin{proposition}[The constant-term kernel in a unipotent group]
\label{prop:BCH}
There is a polynomial-coordinate identification
\[
 \ker\bigl(U(\Bk)\longrightarrow U(k)\bigr)
   =\exp\bigl(\mathfrak u(k)\otimes_k\I_k\bigr).
\]
On the right, multiplication corresponds to the finite
Baker--Campbell--Hausdorff polynomial. The tensor notation means finite
linear combinations of a fixed finite $k$-basis of $\mathfrak u$ with
coefficients in $\I_k$.
For an integral model with unipotent $k$-fiber, its constant-term identity
fiber is the same group inside that faithful representation.
\end{proposition}
\begin{proof}
Exponential and logarithm are polynomials over $k$, so they commute with
$\ct:\Bk\to k$. A Lie-algebra point has zero constant term exactly
when its coordinates in a $k$-basis lie in $\I_k$. The polynomial
exponential has constant term the identity precisely in that case.
The Baker--Campbell--Hausdorff expression truncates because all sufficiently
long brackets vanish, and its remaining brackets again have coefficients
in $\I_k$. For an integral model, these points and their inverses have
identity constant matrix and purely infinite corrections. They satisfy
its defining equations as in \cref{cor:expI}. The converse follows by
embedding into the $k$-fiber and taking the logarithm.
\end{proof}

Nothing in this proposition asserts convergence of a general exponential
series at an infinite surreal argument. Every displayed matrix function
is a finite polynomial.

\subsection{Upper triangular groups}
Put $U_n=\UT_n(\I_k)$, the group of upper unitriangular $n\times n$
matrices whose entries strictly above the diagonal belong to $\I_k$.
For $r\ge1$ let
\[
 F_r=\{g\in U_n:g_{ij}=0\text{ whenever }0<j-i<r\}.
\]
Thus $F_1=U_n$ and $F_n=1$.

\begin{theorem}[Exact commutator filtration]\label{thm:UT}
For all positive integers $r,s$,
\[
 [F_r,F_s]=F_{r+s}.
\]
Consequently the lower central series is $\gamma_r(U_n)=F_r$, the
successive derived groups are $U_n^{(j)}=F_{2^j}$, and for $n\ge2$
\[
 \text{nilpotency class}(U_n)=n-1,\qquad
 \text{derived length}(U_n)=\lceil\log_2n\rceil.
\]
The first superdiagonal induces
\[
                     U_n/[U_n,U_n]\cong(\I_k,+)^{n-1}.
\]
\end{theorem}
\begin{proof}
Write $F_r=1+J_r$, where $J_r$ consists of the corresponding strictly
upper triangular matrices. Matrix multiplication gives
$J_rJ_s\subseteq J_{r+s}$. Finite geometric inversion of $1+X$
then implies $[F_r,F_s]\subseteq F_{r+s}$.

Elementary elimination expresses every member of $F_t$ as a finite
product of $e_{ij}(a)$ with $j-i\ge t$ and $a\in\I_k$: eliminate
successive superdiagonals, noting that products affect only later ones.
For a generator of $F_{r+s}$, choose $h=i+r$, so $h-i\ge r$ and
$j-h\ge s$. By \cref{prop:factor}, write $a=bc$ with $b,c\in\I_k$.
The elementary identity
\[
                   [e_{ih}(b),e_{hj}(c)]=e_{ij}(bc)
\]
places the generator in $[F_r,F_s]$. If $r+s\ge n$, both sides are
trivial. This proves equality in all cases.

Induction yields the two series. The last nontrivial lower-central term
contains $e_{1n}(a)$ for $a\ne0$, and $F_t\ne1$ exactly when $t<n$.
The first superdiagonal is additive under multiplication, is surjective
onto $(\I_k)^{n-1}$, and has kernel $F_2$.
\end{proof}

For any fixed positive surreal exponent $\rho$, coefficient extraction
$f\mapsto[\omega^\rho]f$ is a nonzero additive map $\I_k\to k$.
Composing it with a first-superdiagonal coordinate gives a nontrivial
homomorphism $U_n\to(k,+)$ when $n\ge2$. Thus a purely infinite
unipotent group can be visible to small groups. The opposite-root
commutators in $P_n(k)$ will destroy all such visibility.

\status{The exponential correspondence and unitriangular commutator
calculus are classical. The exact formulas here use the already established
single-product property of $\I_k$. Their purpose is to describe the
nonrigid alternative and to clarify why the elementary-kernel theorem is
not a statement about all purely infinite matrix groups.}

\section{A support-absorption lemma}\label{sec:absorption}

The next argument is where the full surreal exponent class, rather than
an arbitrary Hahn workspace, becomes indispensable. It is a version of
the geometric collision construction in \cite[Section 3]{RepoQuot},
adapted to produce an arbitrary prescribed root parameter.

\begin{lemma}[Absorption by a difference of nearby monomials]
\label{lem:absorb}
Fix $0\ne c\in\I_k$. There is $\varepsilon>0$ such that for any
$0<\beta<\alpha<\varepsilon$,
\begin{equation}\label{eq:absorption}
 \frac{c}{\omega^\alpha-\omega^\beta}
   =\sum_{m\ge0}c\,\omega^{-\alpha-m(\alpha-\beta)}\in\I_k.
\end{equation}
For every set cardinal $\lambda$, the interval $(0,\varepsilon)$
contains a set of $\lambda$ distinct surreal exponents.
\end{lemma}
\begin{proof}
Choose $\delta>0$ strictly below $\supp(c)$ using \cref{lem:gap}, and
put $\varepsilon=\delta/(4\omega)$, with this arithmetic performed in
the surreal \emph{exponent field}. For every positive ordinary integer
$m$, $m\varepsilon<\delta$.

Set $\gamma=\alpha-\beta>0$. By the Hahn geometric identity,
\[
 (\omega^\alpha-\omega^\beta)^{-1}
     =\omega^{-\alpha}(1-\omega^{-\gamma})^{-1}
     =\sum_{m\ge0}\omega^{-\alpha-m\gamma}.
\]
Multiplying this Hahn series by the single series $c$ is legitimate by
the support lemma. Every exponent that can occur in the product has
the form $s-\alpha-m\gamma$, with $s\in\supp(c)$; and
\[
 s-\alpha-m\gamma
    >\delta-(m+1)\varepsilon>0.
\]
Cancellations can remove terms but cannot introduce a nonpositive exponent.
Thus the product belongs to $\I_k$ and proves \eqref{eq:absorption}.

For an ordinal $\xi$, regard $\xi$ as its embedded surreal ordinal and
put $\alpha_\xi=\varepsilon/(2+\xi)$, using surreal addition and
division. These are distinct and lie in $(0,\varepsilon)$. Restricting
to $\xi<\lambda$ gives the required set.
\end{proof}

\begin{remark}[The infinite series is not mapped term by term]
\label{rem:nomapseries}
If $d=\omega^\alpha-\omega^\beta$ and $b$ denotes the entire series
on the right of \eqref{eq:absorption}, the only identity needed below is
$\,db=c$. A target group homomorphism is applied to a finite matrix
commutator involving $b$; it is never applied to the summands of the
series. No continuity or infinite additivity is smuggled into the proof.
\end{remark}

For finite truncations $b_M$ of \eqref{eq:absorption}, the exact formula is
\begin{equation}\label{eq:remainder}
       d b_M=c\bigl(1-\omega^{-(M+1)(\alpha-\beta)}\bigr).
\end{equation}
The remainder is not zero. The finite identity is useful for checking
signs, but replacing $b$ by a truncation would not prove the theorem.

\section{Purely infinite elementary groups}\label{sec:elementary}

\subsection{Normality and perfectness}
Throughout this section $n\ge3$. We use the convention
$[x,y]=xyx^{-1}y^{-1}$. The elementary identities
\begin{align}
 e_{ij}(a)e_{ij}(b)&=e_{ij}(a+b),\label{eq:rootadd}\\
 [e_{ih}(a),e_{hj}(b)]&=e_{ij}(ab)
                    \quad(i,h,j\text{ distinct})\label{eq:rootcomm}
\end{align}
follow by multiplying matrix units. These are classical; see also
\cite[Chapter III, Section 1]{Weibel}.

\begin{proposition}[An abstract idempotent-ideal lemma]\label{prop:normal}
Let $A$ be a commutative unital ring and $I$ an ideal such that every
$a\in I$ is a product of two elements of $I$. Then
\[
 P=\langle e_{ij}(a):a\in I,\ i\ne j\rangle
\]
is perfect and normal in $E_n(A)$.
\end{proposition}
\begin{proof}
For a generator $e_{ij}(a)$, choose a third index $h$ and a factorization
$a=bc$ with $b,c\in I$. Equation \eqref{eq:rootcomm} expresses this
generator as a commutator of two elements of $P$. Hence $P=[P,P]$.

To prove normality, conjugate a generator by $z=e_{kl}(r)$, $r\in A$.
Except in the opposite-root case $(k,l)=(j,i)$, the conjugate is either
unchanged or a product of root elements with parameters $a$ and $\pm ra$.
Those parameters belong to $I$. For the opposite root, write
$e_{ij}(a)=[e_{ih}(b),e_{hj}(c)]$ as above. Direct multiplication gives
\begin{align}
 e_{ji}(r)e_{ih}(b)e_{ji}(-r)
      &=e_{jh}(rb)e_{ih}(b),\label{eq:conjone}\\
 e_{ji}(r)e_{hj}(c)e_{ji}(-r)
      &=e_{hi}(-cr)e_{hj}(c).\label{eq:conjtwo}
\end{align}
Both conjugated factors belong to $P$. Their commutator therefore does
as well. Since the elementary matrices and their inverses generate
$E_n(A)$, $P$ is normal.
\end{proof}

A version with $I^2=I$ and finite sums of products follows by splitting
$e_{ij}(a)$ according to \eqref{eq:rootadd}. We have stated the stronger
single-product hypothesis because \cref{prop:factor} gives it directly.

\begin{proposition}[Exact split kernel]\label{prop:split}
For any unital subring $D\subseteq k$, with $A=D\oplus\I_k$,
\[
 E_n(A)=P_n(k)\rtimes E_n(D),\qquad
 P_n(k)=\ker\bigl(\ct:E_n(A)\to E_n(D)\bigr).
\]
\end{proposition}
\begin{proof}
Normality follows from \cref{prop:factor,prop:normal}. For $a=d+b$ with
$d\in D$ and $b\in\I_k$,
$e_{ij}(a)=e_{ij}(b)e_{ij}(d)$. Normality allows the purely infinite
factors in any finite word to be moved to the left. Hence every $g$ can
be written $g=pq$, with $p\in P_n(k)$ and $q\in E_n(D)$.
Applying $\ct$ gives $\ct(g)=q$. Thus the kernel is exactly $P_n(k)$,
and its intersection with $E_n(D)$ is trivial. The inclusion of constant
matrices is a section, proving the semidirect product.
\end{proof}

For $D=\Z$ or $\Z[i]$, Euclidean row reduction proves
$E_n(D)=\SL_n(D)$. Indeed, the entries of the first column of a
determinant-one matrix generate the unit ideal. Euclidean elementary
operations reduce that column to a unit times the first basis vector;
a determinant-one unit diagonal block is elementary, so the unit can be
removed. Clear the first row and induct on the remaining block. Over
$\Z[i]$ the Euclidean function is the usual Gaussian norm. These are
ordinary finite matrix arguments, unrelated to a Euclidean algorithm
in $\Oz$, which we do not assume.

\subsection{No small abstract representation of the kernel}
\begin{theorem}[The purely infinite kernel is invisible to sets]
\label{thm:invisible}
Every class group homomorphism $\phi:P_n(k)\to H$ into a set-sized
group $H$ is trivial.
\end{theorem}
\begin{proof}
It suffices to prove $\phi(e_{ij}(c))=1_H$ for each generator. The case
$c=0$ is immediate; fix $c\ne0$ and a third index $h$.
Apply \cref{lem:absorb} to $c$, obtaining $\varepsilon$.
Choose a set $T\subset(0,\varepsilon)$ with $|T|>|H|$.
The set of elements
\[
                        e_{ih}(\omega^\alpha),\qquad \alpha\in T,
\]
cannot have distinct images in $H$. Choose distinct $\alpha,\beta\in T$
with equal images, and interchange their names so that $\alpha>\beta$.
By \eqref{eq:rootadd},
\[
                 \phi\bigl(e_{ih}(\omega^\alpha-\omega^\beta)\bigr)=1_H.
\]
Put $d=\omega^\alpha-\omega^\beta$. By \cref{lem:absorb},
$b=c/d$ belongs to $\I_k$. Equation \eqref{eq:rootcomm} now gives
\[
 \phi(e_{ij}(c))
     =\phi([e_{ih}(d),e_{hj}(b)])
     =[1_H,\phi(e_{hj}(b))]=1_H.
\]
Every generator is killed, so every finite word is killed.
\end{proof}

This proof does not first construct a ring map from $A$ to $H$, or from
$A$ to an endomorphism ring. It applies to arbitrary abstract group maps
on $P_n(k)$ itself. It also works for $k=\C$ without an ordering on the
coefficients: positivity is required only of exponents.

\begin{proof}[Proof of \cref{main:elementary}]
Perfectness, normality, and the split kernel follow from
\cref{prop:normal,prop:split}, the single-product property, and elementary
generation over $D=\Z$ or $\Z[i]$. Nontriviality follows from
$e_{12}(\omega)\ne1_n$. Given a homomorphism $\psi:E_n(A)\to H$,
its restriction to $P_n(k)$ is trivial by \cref{thm:invisible}.
It therefore factors uniquely through $\ct$ in the split decomposition.
Conversely any homomorphism $\SL_n(D)\to H$ composes with $\ct$.
\end{proof}

\begin{corollary}[Actions and linear representations]\label{cor:actions}
Every action of $P_n(k)$ on a set is trivial. Every action of $E_n(A_{D,k})$
on a set, for $D=\Z$ or $\Z[i]$, factors uniquely through its constant
$\SL_n(D)$ quotient. In particular, every linear representation on a
set-sized vector space over any set-sized field kills $P_n(k)$, regardless
of the dimension of the vector space.
\end{corollary}
\begin{proof}
An action on a set $X$ is a homomorphism to the set-sized group $\Sym(X)$.
For linear representations use the set-sized group of invertible linear
maps of the specified vector space.
\end{proof}

\begin{corollary}[Other constant coefficient rings]\label{cor:otherD}
For any set-sized unital subring $D\subseteq k$, the universal set-sized
abstract group quotient of $E_n(D\oplus\I_k)$ is $E_n(D)$.
In particular, for $D=k$, it is $\SL_n(k)$.
\end{corollary}
\begin{proof}
The proof uses only \cref{prop:split,thm:invisible}. The group $E_n(D)$
is set-sized because $D$ is. Over the field $k$, ordinary elimination gives
$E_n(k)=\SL_n(k)$.
\end{proof}

Perfectness alone would only kill homomorphisms to abelian groups.
\Cref{thm:invisible} kills homomorphisms to \emph{all} set-sized groups,
including nonabelian and infinite groups. It does not make $P_n(k)$
trivial: its defining faithful matrix representation has a proper-class
coefficient field, hence does not have a set-sized target.

\section{A separating proper-class congruence tower}\label{sec:tower}

A group with no small images might be mistaken for a simple group. The
kernels just constructed provide a particularly explicit counterexample
to that inference.

Fix $A=\Oz$ or $\GOz$, and write $P=P_n(k)$. For each positive
surreal exponent $a$, define
\begin{equation}\label{eq:Na}
 N_a=\{g\in P:g-1_n\in M_n(\omega^a A)\}.
\end{equation}
Equivalently, $N_a$ is the kernel on $P$ of matrix reduction modulo the
principal ideal $\omega^a A$. This description can be used just as a
congruence relation; no set-sized quotient ring is assumed.

\begin{theorem}[A non-simple invisible group with a large residual tower]
\label{thm:tower}
For every $a>0$, $N_a$ is a nontrivial proper normal subgroup of $P$.
If $0<a<b$, then $N_b\subsetneq N_a$. Moreover,
\begin{equation}\label{eq:towerintersection}
 \bigcap_{a\in\No^{>0}}N_a=1,
 \qquad
 \bigcap_{a\in S}N_a\ne1
 \quad\text{for every set }S\subset\No^{>0}.
\end{equation}
Every nontrivial quotient $P/N_a$ has no nontrivial homomorphism into a
set-sized group and cannot itself be set-sized.
\end{theorem}
\begin{proof}
For a two-sided ideal $J$ in a commutative ring, the matrices congruent
to the identity modulo $J$ form a normal subgroup of $\GL_n(A)$:
products, inverses, and conjugation preserve this congruence. Intersecting
with $P$ proves normality of $N_a$.

The element $e_{12}(\omega^a)$ belongs to $N_a$. But
$e_{12}(\omega^{a/2})$ does not, because its nonzero off-diagonal entry
would require a quotient $\omega^{-a/2}\in A$, which is impossible.
Thus $N_a$ is nontrivial and proper.

For $b>a$, $\omega^b A\subseteq\omega^a A$, since
$\omega^{b-a}\in A$. The element $e_{12}(\omega^a)$ witnesses strictness:
$\omega^a/\omega^b$ has negative exponent and is not in $A$.

If $g\ne1_n$, choose a nonzero entry $f$ of $g-1_n$. For any
$a>\bdeg(f)$, the quotient $\omega^{-a}f$ has negative leading exponent
and cannot lie in $A$. Therefore $g\notin N_a$, proving the first
intersection formula. For a set $S$ of positive exponents, choose
$b>a$ for every $a\in S$ using \cref{lem:gap}. Then
$e_{12}(\omega^b)$ is a nonidentity element of every $N_a$, proving
the second formula, including the empty-set case.

Finally, a homomorphism from $P/N_a$ to a set-sized group composes with
the quotient map to a homomorphism from $P$, hence is trivial. If the
nontrivial quotient had a set-sized carrier, its identity homomorphism
would contradict this conclusion.
\end{proof}

Each $P/N_a$ is also perfect, being a quotient of a perfect group. Thus
one obtains explicit nontrivial perfect class groups with no nontrivial
small images at every level of a strictly descending congruence tower.
The tower separates individual matrices only when all surreal scales are
allowed. This is a concrete set/class distinction, rather than a claim
about a numerical ``size'' assigned to $P$.

\section{Sharp boundaries of the argument}\label{sec:boundaries}

\subsection{A fixed Hahn workspace is not the full class}
Let $\Gamma$ be a nonzero set-sized divisible ordered abelian group, and
use large-monomial notation $X^\gamma$. Form
\[
 \I_k(\Gamma)=\{f:\supp f\subseteq\Gamma^{>0}\},\qquad
 A_{D,k}(\Gamma)=D\oplus\I_k(\Gamma)
\]
inside the Hahn field with set-sized reverse-well-ordered supports.
These rings are sets.

The geometric dichotomy persists for their coefficient-enlarged versions:
if $k=\R$, the divisible-group real Hahn field is real closed; if $k=\C$,
it is algebraically closed. Unit rigidity and bounded-coordinate rigidity
are unchanged, and a positive exponent gives a nonconstant parameter.
Thus the proofs of \cref{main:real,main:complex} apply with the corresponding
Hahn-cone rings in place of $\BR,\BC$ and with integral constant rings
$D$ as before. The complex Hahn closedness assertion follows from
\cite[Corollary 4]{Poonen}. The real Hahn field is ordered, and its
complexification is that algebraically closed complex Hahn field; hence
it is real closed.

In contrast, universal invisibility cannot persist in a nontrivial
set-sized workspace: its purely infinite elementary group is itself a
set-sized nontrivial group, and its identity homomorphism is faithful.
The full-class proof needs arbitrarily many scales below a support gap.

\begin{proposition}[Local cardinal form]\label{prop:local}
Fix $c\ne0$ in $\I_k(\Gamma)$ and suppose that $0<\delta<s$ for all
$s\in\supp(c)$. Let $T\subseteq\Gamma^{>0}$ satisfy
\[
                   m\alpha<\delta
       \quad(\alpha\in T,\ m\in\N_{>0}).
\]
For a homomorphism from the purely infinite elementary group of rank
$n\ge3$ over this workspace to a group $H$ with $|T|>|H|$, every root
element $e_{ij}(c)$ maps to the identity.
\end{proposition}
\begin{proof}
The collision proof of \cref{thm:invisible} applies verbatim. For the two
selected exponents $\alpha>\beta$,
$\alpha+m(\alpha-\beta)\le(m+1)\alpha<\delta$, so the quotient
$c/(X^\alpha-X^\beta)$ lies in the same full Hahn workspace and has
strictly positive support. No statement about normality or perfectness
in this workspace is needed.
\end{proof}

This is a conditional lower bound on target size, not an assertion that
every exponent group supplies the required set $T$. For an Archimedean
exponent group no positive $\alpha$ satisfies all $m\alpha<\delta$.

\subsection{Finite support changes the answer}
Let $\I_k^{\mathrm{fin}}$ consist only of the purely infinite normal forms
with finite support, and put $A^{\mathrm{fin}}=D+\I_k^{\mathrm{fin}}$.
The exponent class is still all of $\No$.

\begin{proposition}[A visible finite-support kernel]\label{prop:finsupport}
The finite sum of coefficients defines a ring homomorphism
\[
 \operatorname{aug}:A^{\mathrm{fin}}\longrightarrow k,
 \qquad \sum_s a_s\omega^s\longmapsto\sum_s a_s.
\]
For $n\ge3$, it maps the group generated by the purely infinite elementary
matrices \emph{onto} the set-sized group $\SL_n(k)$.
\end{proposition}
\begin{proof}
Both sums and products involve only finite supports, so expanding a product
and summing its coefficients proves multiplicativity. For every $r\in k$
and every positive exponent $a$, the parameter $r\omega^a$ is purely
infinite and maps to $r$. Thus the induced matrix map contains every
$e_{ij}(r)$ in its image. These generate $\SL_n(k)$.
\end{proof}

The finite-support ideal still has the single-product property, so the
corresponding elementary kernel is perfect and normal. Nevertheless it
has a large, ordinary image. Infinite Hahn supports are therefore not
cosmetic in \cref{thm:invisible}: they supply the exact divisor in
\eqref{eq:absorption}, which finite truncations cannot replace.

\subsection{Rank two and the full special linear group}
The proof of \cref{thm:invisible} needs a third matrix index in order to
convert a killed additive root parameter into an arbitrary multiplicative
one. It makes no assertion about rank $2$. Similarly, the exact theorem
concerns $E_n(A)$; it does not assert $E_n(A)=\SL_n(A)$ or silently use
a stable-range theorem for an omnific ring. These are separate issues.

The two natural next questions are whether the universal small quotient
of $\SL_n(\Oz)$ is $\SL_n(\Z)$, and whether an analogue of kernel
invisibility holds in rank $2$. These are questions left by the present
arguments, not claims that the literature has designated them as open
problems. For the first, non-elementary matrices and relative algebraic
$K$-theory are the additional obstruction; for the second, the missing
third-root commutator must be replaced by another mechanism.

\section{Synthesis and proof dependencies}\label{sec:synthesis}

The two main phenomena should not be conflated. The geometric theorem is
a finite-dimensional statement controlled by real or complex algebraic
group structure. It holds in suitably chosen set-sized Hahn fields as well
as in the full surreal class. The abstract quotient theorem is a genuinely
class-sized phenomenon and depends on both arbitrary small exponent scales
and infinitely supported source elements.

The logical dependencies can be summarized without relying on any
unproved assertion of a companion draft:
\[
 \begin{gathered}
 \text{normal-form units}\ \Longrightarrow\ \text{torus freezing},\\
 \text{bounded real coordinates + real-closed transfer}
       \ \Longrightarrow\ \text{bounded-fiber rigidity},\\
 \text{classical group structure + polynomial unipotent curves}
       \ \Longrightarrow\ \text{the two group dichotomies};
 \end{gathered}
\]
\[
 \begin{gathered}
 \text{set-sized support gap}\ \Longrightarrow\ \text{single-product divisibility},\\
 \text{Hahn geometric inverse + arbitrarily many small scales}
       \ \Longrightarrow\ \text{support absorption},\\
 \text{support absorption + a third elementary index}
       \ \Longrightarrow\ \text{no set-sized kernel representations}.
 \end{gathered}
\]
The normality and perfectness proofs use only finite matrix identities and
single-product divisibility. The universal quotient then follows from a
split constant-term map. The congruence tower uses only principal ideals,
leading exponents, and upper bounds for sets of exponents.

In particular, neither an unproved omnific factorization conjecture nor
a global analytic theory on $\No[i]$ is needed. Published factorization
work such as \cite{LM} supplies important context but is not a hidden
hypothesis for any of the three main theorems.

\appendix
\clearpage
\section{Repository comparison and provenance}\label{app:audit}

The repository snapshot used for comparison is
\begin{center}\pin.\end{center}
The inspection date is 23 September 2026. The repository was read through
its connected GitHub interface; it was not modified, cloned successfully,
or built in this session. No proof-assistant verification of the new
statements is claimed.

The catalogue describes 51 AI-assisted research reports and separates
mathematical drafts from checked formalization \cite{RepoCatalogue}.
The comparison covered that catalogue and the opening 280 source lines
of the omnific Diophantine article \cite{RepoDio}. For the universal
ring-quotient article it covered the main statements and the subsequently
returned sections on class conventions, common divisibility, and
geometric-series collisions \cite{RepoQuot}. This was targeted source
inspection, not an audit of every report or historical manuscript.

Several boundaries are particularly important. The earlier Diophantine
article already contains constant-product rigidity, Pell and norm
rigidity, a quadratic-level dichotomy, and constructions of infinite
families. We do not reclassify these as new results. The earlier quotient
article already proves that every unital map from $\Oz$ into a set-sized
ring factors through $\Z$, and develops module and Gaussian analogues.
Its support-gap and geometric-series collision arguments are acknowledged
as the starting point for \cref{sec:absorption}.

The catalogue also describes a universal set-sized quotient for
$\SO(3,\No)$ in the Euclidean three-space report. Only that description,
not a full source proof audit of the report, was used here. Thus the
\emph{idea} of a universal small group quotient is not new to the
repository. Our source group is instead the elementary group over the
omnific ring, and the invisible subgroup is generated by purely
\emph{infinite} root parameters. The standalone assertion about all
abstract maps on $P_n(k)$ is the distinctive strengthening.

\begin{center}
\small
\begin{tabular}{P{0.35\textwidth}P{0.56\textwidth}}
\toprule
Item & Attribution and status \\
\midrule
Normal forms, Hahn products, real closedness & Classical; cited background \\
Constant term, units, support gaps and geometric collision & Background and prior repository material; needed versions reproved \\
Cartan theorem, torus quotient, nilpotent exponential & Classical; explicit imported inputs \\
Real and Gaussian group dichotomies & Proposed combined contributions, with proofs \\
Abstract invisibility of $P_n(k)$ and elementary universal quotient & Proposed group-level extension, with a direct proof on the kernel \\
Separating proper-class congruence tower & Consequence proved here; no simplicity claim \\
Unitriangular filtration and finite-support contrast & Explicit supporting deductions; no separate priority claim \\
\bottomrule
\end{tabular}
\end{center}

The literature search included the author-hosted algebraic-group and
$K$-theory sources in the bibliography and queries combining ``omnific''
with ``unipotent'', ``algebraic groups'', and ``elementary matrices''.
No source giving the exact new combined statements was located. This
search cannot exclude an earlier result in another terminology or an
unindexed document. Independent mathematical and historical review would
still be required before claiming a publishable original theorem.

\clearpage
\section{Finite verification and reproducibility}\label{app:verification}

The accompanying \texttt{verification.py} performs exact calculations
with SymPy, not floating-point experiments. The included run used Python
3.13.5 and SymPy 1.14.0 and passed 1,108 checks. One matrix identity counts
as one check after every entry has been tested for exact equality.

\begin{center}
\begin{tabular}{lr}
\toprule
Check family & Exact checks passed \\
\midrule
Elementary addition, commutators, and conjugation & 908 \\
The two orthogonal one-parameter examples & 18 \\
Unitriangular filtration instances & 139 \\
Geometric remainders and finite-support augmentation & 43 \\
\midrule
Total & 1,108 \\
\bottomrule
\end{tabular}
\end{center}


The program checks the additive root identity, the third-index commutator,
and both opposite-root conjugation formulas for every relevant ordered
pair or triple of indices in ranks $3$ through $6$. It checks the combined
normality identity as well. For each of \eqref{eq:complexexample} and
\eqref{eq:realexample}, it verifies nilpotence, the skew relation,
orthogonality, determinant one, the addition and inverse formulas,
truncated logarithmic recovery, and a coordinate recovering the parameter.

It also checks exact integer instances of the unitriangular filtration,
the finite geometric remainder \eqref{eq:remainder} in a formal variable,
and multiplicativity of finite-support augmentation. The integer
filtration instances illustrate ring identities; they are not claimed to
be representations of arbitrary elements of $\I_k$.

The program does \emph{not} verify the Hahn support lemma, the existence
of surreal gaps below arbitrary sets, the proper-class collision argument,
Cartan's theorem, or the universal quantifiers in the main results. Those
are addressed by the written proofs and cited classical inputs. A
successful symbolic run and a successful PDF build do not establish
those facts independently.

To reproduce the files, run the following commands from the extracted
archive directory:
\begin{verbatim}
python verification.py --output verification-results.txt
pdflatex -interaction=nonstopmode -halt-on-error article.tex
pdflatex -interaction=nonstopmode -halt-on-error article.tex
pdflatex -interaction=nonstopmode -halt-on-error article.tex
\end{verbatim}
The LaTeX source contains its own bibliography and needs no external
bibliographic database. Standard TeX Live packages suffice. SymPy is
needed only for the optional finite checks, not for compiling the article.

\clearpage
\phantomsection
\addcontentsline{toc}{section}{References}
\begin{thebibliography}{99}
\small\raggedright
\setlength{\itemsep}{5pt}

\bibitem{Conway}
J. H. Conway, \emph{On Numbers and Games}, second edition,
A K Peters, 2001.

\bibitem{Gonshor}
H. Gonshor, \emph{An Introduction to the Theory of Surreal Numbers},
London Mathematical Society Lecture Note Series 110,
Cambridge University Press, 1986.

\bibitem{Berarducci}
A. Berarducci, \emph{Surreal numbers, exponentiation and derivations},
2020, arXiv:2008.06878.
\url{https://arxiv.org/abs/2008.06878}.

\bibitem{LM}
S. L'Innocente and V. Mantova,
\emph{A factorisation theory for generalised power series and omnific integers},
arXiv:1710.07304v5, 22 January 2024.
\url{https://arxiv.org/abs/1710.07304v5}.

\bibitem{MilneAG}
J. S. Milne, \emph{Algebraic Groups: The Theory of Group Schemes of
Finite Type over a Field}, Cambridge Studies in Advanced Mathematics 170,
Cambridge University Press, 2017.
Author-hosted text: \url{https://www.jmilne.org/math/Books/iAG2017.pdf}.

\bibitem{MilneSV}
J. S. Milne, \emph{Introduction to Shimura Varieties},
Clay Mathematics Proceedings 4 (2005), 265--378;
author's revised notes dated 16 September 2017.
\url{https://www.jmilne.org/math/xnotes/svi.pdf}.

\bibitem{MilneRG}
J. S. Milne, \emph{Reductive Groups}, course notes, version 2.00,
10 March 2018.
\url{https://www.jmilne.org/math/CourseNotes/RG.pdf}.

\bibitem{Marker}
D. Marker, \emph{Model Theory: An Introduction},
Graduate Texts in Mathematics 217, Springer, 2002.

\bibitem{Weibel}
C. A. Weibel, \emph{The K-book: An Introduction to Algebraic K-theory},
Graduate Studies in Mathematics 145, American Mathematical Society, 2013.
Author-hosted Chapter III:
\url{https://sites.math.rutgers.edu/~weibel/Kbook/Kbook.III.pdf}.

\bibitem{Poonen}
B. Poonen, \emph{Maximally complete fields},
L'Enseignement Math\'ematique (2) 39 (1993), 87--106.
\url{https://math.mit.edu/~poonen/papers/amsval.pdf}.

\bibitem{RepoCatalogue}
\texttt{VladimirReshetnikov/Surreal},
\emph{The surreal and surcomplex reports}, documentation catalogue,
snapshot \pin, inspected 23 September 2026.
\url{https://github.com/VladimirReshetnikov/Surreal/blob/220784a4766b9dd9ab3deb92f9f28ddfdd866187/docs/README.md}.

\bibitem{RepoDio}
\texttt{VladimirReshetnikov/Surreal},
\emph{Omnific Integers and Omnific--Diophantine Geometry:
Retractions, rigidity, definability, and infinite families},
AI-assisted research draft dated 22 September 2026,
source at the snapshot recorded in \cref{app:audit}.
\url{https://github.com/VladimirReshetnikov/Surreal/blob/220784a4766b9dd9ab3deb92f9f28ddfdd866187/docs/surreal/omnific-diophantine-geometry/article.tex}.

\bibitem{RepoQuot}
\texttt{VladimirReshetnikov/Surreal},
\emph{The Universal Set-Sized Quotient of the Omnific Integers:
Cardinal thresholds, surcomplex representations, and homological obstructions},
AI-assisted research draft dated 22 September 2026,
source at the snapshot recorded in \cref{app:audit}.
\url{https://github.com/VladimirReshetnikov/Surreal/blob/220784a4766b9dd9ab3deb92f9f28ddfdd866187/docs/surreal/set-sized-quotients-of-omnific-integers/article.tex}.

\end{thebibliography}
\end{document}
